\documentclass[aspectratio=169]{beamer}

\usetheme{Madrid}
\usepackage[utf8]{inputenc}
\usepackage[T1]{fontenc}
\usepackage[french]{babel}
\usepackage{amsmath,amssymb,amsfonts}
\usepackage{bm}

\title{Machines à vecteurs de support (SVM)}
\subtitle{Cas linéairement séparable -- marge rigide}
\author{D'après le chapitre sur les SVM}
\date{}

\begin{document}

%----------------- Titre -----------------%
\begin{frame}
  \titlepage
\end{frame}

%----------------- Plan -----------------%
\begin{frame}{Plan}
  \tableofcontents
\end{frame}

\section{Problème de classification linéairement séparable}

%----------------- Données et séparabilité -----------------%
\begin{frame}{Problème de classification binaire}
  On considère un problème de classification binaire avec un jeu de données :
  \[
    D = \{(\bm{x}_i, y_i)\}_{i=1}^n,
  \]
  où
  \[
    \bm{x}_i \in \mathbb{R}^p, \qquad y_i \in \{-1, +1\}.
  \]

  \vspace{0.5cm}
  \textbf{Idée :} trouver un \textbf{classifieur linéaire}
  \[
    f(\bm{x}) = \langle \bm{w}, \bm{x} \rangle + b
  \]
  tel que la frontière de décision
  \[
    H = \{\bm{x} \in \mathbb{R}^p \mid \langle \bm{w}, \bm{x} \rangle + b = 0\}
  \]
  sépare les deux classes.
\end{frame}

\begin{frame}{Définition : séparabilité linéaire}
  \textbf{Définition.} Le jeu de données $D$ est dit \emph{linéairement séparable} s'il existe un hyperplan
  \[
    H = \{\bm{x} \in \mathbb{R}^p \mid \langle \bm{w}, \bm{x} \rangle + b = 0\}
  \]
  tel que :
  \[
    \begin{cases}
      \langle \bm{w}, \bm{x}_i \rangle + b > 0 & \text{si } y_i = +1,\\[0.2cm]
      \langle \bm{w}, \bm{x}_i \rangle + b < 0 & \text{si } y_i = -1.
    \end{cases}
  \]

  \vspace{0.3cm}
  Cela revient à demander
  \[
    y_i \bigl( \langle \bm{w}, \bm{x}_i \rangle + b \bigr) > 0 \quad \text{pour tout } i.
  \]

  \vspace{0.3cm}
  Dans ce cas, il existe \textbf{une infinité d'hyperplans séparateurs} qui ne font aucune erreur
  de classification : il faut choisir le ``meilleur''.
\end{frame}

\section{Marge d'un hyperplan séparateur}

\begin{frame}{Marge et hyperplans parallèles}
  \textbf{Idée :} parmi tous les hyperplans séparateurs, on choisit celui qui maximise la \emph{marge}.

  \vspace{0.3cm}
  On considère :
  \[
    H : \langle \bm{w}, \bm{x} \rangle + b = 0,
  \]
  ainsi que deux hyperplans parallèles :
  \[
    H_+ : \langle \bm{w}, \bm{x} \rangle + b = 1, \qquad
    H_- : \langle \bm{w}, \bm{x} \rangle + b = -1.
  \]

  \vspace{0.3cm}
  Les contraintes de séparation deviennent alors :
  \[
    y_i \bigl( \langle \bm{w}, \bm{x}_i \rangle + b \bigr) \geq 1, \quad i=1,\dots,n.
  \]

  \vspace{0.3cm}
  La \textbf{marge} $\gamma$ est la distance entre $H$ et $H_+$ (ou $H_-$) :
  \[
    \gamma = \frac{1}{\|\bm{w}\|_2}.
  \]
\end{frame}

\begin{frame}{Vecteurs de support et zone d'indécision}
  \begin{itemize}
    \item Les observations qui vérifient
      \[
        y_i \bigl( \langle \bm{w}, \bm{x}_i \rangle + b \bigr) = 1
      \]
      sont situées exactement sur $H_+$ ou $H_-$.
    \item Ces points s'appellent les \textbf{vecteurs de support} : ils ``soutiennent'' la solution.
    \item La région comprise entre $H_-$ et $H_+$ s'appelle la \textbf{zone d'indécision} : elle ne contient
          aucune observation dans le cas séparé.
    \item Seuls les vecteurs de support déterminent la position optimale de l'hyperplan.
  \end{itemize}
\end{frame}

\section{Formulation primale de la SVM à marge rigide}

\begin{frame}{Problème d'optimisation primale}
  L'objectif de la SVM à marge rigide est de \textbf{maximiser la marge} $\gamma = 1/\|\bm{w}\|_2$,
  sous les contraintes de séparation.

  \vspace{0.3cm}
  Maximiser $1/\|\bm{w}\|_2$ revient à minimiser $\tfrac{1}{2}\|\bm{w}\|_2^2$.
  
  \vspace{0.3cm}
  \textbf{Problème primal :}
  \[
    \begin{aligned}
      \min_{\bm{w} \in \mathbb{R}^p,\; b \in \mathbb{R}} \quad & \frac{1}{2} \|\bm{w}\|_2^2 \\
      \text{sous contraintes} \quad & y_i \bigl( \langle \bm{w}, \bm{x}_i \rangle + b \bigr) \geq 1,
      \quad i = 1,\dots,n.
    \end{aligned}
  \]

  \vspace{0.3cm}
  Une solution optimale $(\bm{w}^\ast, b^\ast)$ définit la fonction de décision
  \[
    f(\bm{x}) = \langle \bm{w}^\ast, \bm{x} \rangle + b^\ast.
  \]
\end{frame}

\section{Formulation duale et vecteurs de support}

\begin{frame}{Formulation duale : idée générale}
  Le problème primal est convexe avec des contraintes linéaires.  
  On introduit les multiplicateurs de Lagrange $\alpha_i \geq 0$, un par contrainte.

  \vspace{0.3cm}
  Le lagrangien est :
  \[
    \mathcal{L}(\bm{w}, b, \bm{\alpha}) =
      \frac{1}{2} \|\bm{w}\|_2^2
      - \sum_{i=1}^n \alpha_i
      \Bigl( y_i (\langle \bm{w}, \bm{x}_i \rangle + b) - 1 \Bigr).
  \]

  En minimisant $\mathcal{L}$ par rapport à $\bm{w}$ et $b$, puis en maximisant par rapport à $\bm{\alpha}$,
  on obtient le \textbf{problème dual}.
\end{frame}

\begin{frame}{Formulation duale explicite}
  La formulation duale de la SVM à marge rigide s'écrit :
  \[
    \begin{aligned}
      \max_{\bm{\alpha} \in \mathbb{R}^n} \quad
      & \sum_{i=1}^n \alpha_i - \frac{1}{2}
        \sum_{i=1}^n \sum_{j=1}^n
        \alpha_i \alpha_j y_i y_j \langle \bm{x}_i, \bm{x}_j \rangle \\
      \text{sous contraintes} \quad
      & \sum_{i=1}^n \alpha_i y_i = 0, \\
      & \alpha_i \geq 0,\quad i=1,\dots,n.
    \end{aligned}
  \]

  \vspace{0.4cm}
  On en déduit ensuite :
  \[
    \bm{w}^\ast = \sum_{i=1}^n \alpha_i^\ast y_i \bm{x}_i,
  \]
  où $\bm{\alpha}^\ast$ est une solution optimale du dual.
\end{frame}

\begin{frame}{Fonction de décision et vecteurs de support}
  La fonction de décision peut se réécrire uniquement en fonction des données d'entraînement :
  \[
    f(\bm{x}) = \sum_{i=1}^n \alpha_i^\ast y_i \langle \bm{x}_i, \bm{x} \rangle + b^\ast.
  \]

  \vspace{0.3cm}
  \textbf{Vecteurs de support :}
  \begin{itemize}
    \item Les points tels que $\alpha_i^\ast > 0$ sont les vecteurs de support.
    \item Ils sont situés sur les hyperplans $H_+$ ou $H_-$.
    \item Les points avec $\alpha_i^\ast = 0$ n'interviennent pas dans la somme.
  \end{itemize}

  \vspace{0.3cm}
  \textbf{Conclusion :} la solution ne dépend que d'un sous-ensemble des données, portant toute
  l'information géométrique utile.
\end{frame}

\section{Résumé}

\begin{frame}{Résumé}
  \begin{itemize}
    \item En cas de séparabilité linéaire, la SVM à marge rigide cherche l'hyperplan
          maximisant la marge entre les deux classes.
    \item Le problème s'écrit comme une minimisation de $\frac{1}{2}\|\bm{w}\|_2^2$ sous contraintes
          linéaires $y_i(\langle \bm{w}, \bm{x}_i \rangle + b) \geq 1$.
    \item La formulation duale fait intervenir uniquement des produits scalaires
          $\langle \bm{x}_i, \bm{x}_j \rangle$ et les multiplicateurs de Lagrange $\alpha_i$.
    \item La solution finale dépend uniquement des \textbf{vecteurs de support}
          (les points pour lesquels $\alpha_i^\ast > 0$).
  \end{itemize}
\end{frame}

\end{document}